\documentclass[svgnames,smaller,10pt]{beamer}

\usetheme{metropolis}

\usepackage[utf8]{inputenc}
\usepackage[french]{babel}
\usepackage{hyperref}
\usepackage[T1]{fontenc}
\usepackage{geometry}
\geometry{
    left=2.5mm,
    right=2.5mm,
    top=0mm,
    bottom=0mm
}
\usepackage{amssymb} % pour \blackdiamond
\usepackage{tikz}
\usepackage{graphicx,animate}

\newcommand{\replacementchar}{%
  \tikz[baseline=-0.6ex]{
    %\node[inner sep=0pt] (d) {\Large$\blackdiamond$};
    \node[inner sep=0pt] (d) {\Large$\blacksquare$};
    \node[white, font=\scriptsize\bfseries] at (d.center) {?};
  }%
}
\usepackage{tabularx}
%\usepackage[style=alphabetic,citestyle=authoryear]{biblatex}
%\addbibresource{bibliographie.bib}

\usepackage{xspace}
\newcommand{\themename}{\textbf{\textsc{metropolis}}\xspace}
\title{L2S4: Édition numérique\\
TD 5: Publier sur le Web\\
\small
Renouveler les protocoles éditoriaux
}

\date{}
\author{Marine Tiger\\  \quad {marine.tiger@sorbonne-universite.fr\\}}

\institute{Littératures françaises et comparée ED19, CELLF, Sorbonne Université}

\begin{document}
\maketitle
\setbeamertemplate{section in toc}[sections numbered]

\section{Introduction}

\begin{frame}[c]{Introduction}
\begin{block}{De l'édition à l'\textit{éditorialisation} : comprendre ce que les technologies de l'Internet et du Web font à l'édition (et vice versa).}
    
\end{block}

\begin{itemize}
    \item Les modes de production (technologies? Langages informatiques?)?
    \item Les modes de diffusion?
    \item Les modes de légitimation?
    
\end{itemize}

\end{frame}

\section{Le protocole éditorial}
\begin{frame}[c]{Définition}
\begin{block}{}
    \textbf{Protocole éditorial} = ensemble de règles \textit{explicites} et d'usages en vigueur au sein d'une structure éditoriale.\\
    \textrightarrow guider le processus de conception d'une publication.
\end{block}
\begin{block}{}
    Se présente souvent sous la forme d'une \textbf{série d'instructions précises et détaillées.}\\
\vspace{1em}
    L'ensemble des acteurs de la chaîne éditoriale sont soumis \textbf{au respect du protocole éditorial} (auteurs, éditeurs, préparateurs de copie, designer, directeurs artistiques, imprimeurs...).
\end{block}

\end{frame}

\begin{frame}{Un protocole négocié entre les acteurs et les outils de l'édition}
\begin{columns}[T,onlytextwidth]
\column{0.50\linewidth}
\small
\textrightarrow Détermine un \textbf{choix d'outillage} et donc des choix technologiques.\\
\vspace{1em}
\textrightarrow Vont aussi influencer, en retour, le protocole.
\column{0.50\linewidth}
\includegraphics[width=.90\columnwidth]{05_Axe1_protocoles/images/protocole-hn-auteurs.png}
\end{columns}
    
\end{frame}

\section{Les protocoles informatiques}
\begin{frame}{Définition}
\begin{block}{}
    Ensemble de règles qui permettent à deux équipements informatiques (ordinateurs, smartphones, serveurs) d'échanger des données de manière compréhensible et ordonnée.
\end{block}

\begin{block}{Les protocoles de l'Internet et du Web}
    \begin{itemize}
        \item TCP (\textit{Transmission Control Protocol})
        \item IP (\textit{Internet Protocol})
        \item HTTP (\textit{Hypertext Transfer Protocol})
    \end{itemize}
\end{block}
    
\end{frame}

\begin{frame}{Le protocole HTTP}
    \begin{block}{}
        \centering
        {\small \textit{Le web, ou plus précisément le World Wide Web, est un système documentaire construit sur Internet dans lequel les documents, nommés hypertextes ou pages web, sont reliés les uns aux autres par des hyperliens.Ces documents sont affichés dans des navigateurs qui permettent, grâce aux hyperliens, de naviguer d’une page à une autre. Les données qui sont échangées sur le web le sont avec le protocole HTTP, HyperText Transfert Protocol.}}\\
        \textbf{Jean-Philippe Magué, "Les protocoles d’Internet et du web", Pratiques de l'édition numérique, PUM, 2024}
    \end{block}

    \begin{block}{}
        \textrightarrow{} Fonctionne sur le principe de requête et réponse.
    \end{block}
\end{frame}

\begin{frame}{Apprendre à lire une url}
\begin{block}{}
    \textbf{URL} = protocole, adresse d'un serveur, emplacement et le nom d'une ressource \\
    \textbf{Ex}: \url{https://fr.wikipedia.org/wiki/Wikipédia:Pastiches/Une_femme}
\end{block}
\begin{block}{}
    \begin{itemize}
        \item https:// = le protocole
        \item fr.wikipedia.org = le nom de domaine
        \item /wiki/Wikipédia:Pastiches/Une\_femme = le chemin
    \end{itemize}
\end{block}

\begin{block}{Structure interne de la page Web}
    \begin{itemize}
        \item wiki = le nom de l'outil d'édition
        \item wikipédia: = espace de noms. Signifie que nous sommes dans les pages communautaires ou d'administration du site.
        \item Pastiches = sous-catégorie ou dossier regroupant des textes humoristiques
        \item Une\_femme = titre de la page.
    \end{itemize}
\end{block}
    
\end{frame}

\begin{frame}{Les protocoles éditoriaux de l'Internet : une émanation de la contre-culture américaine}

\begin{block}{Fondements d'une contre-culture basée sur un idéal d'accessibilité des ressources et de liberté d'échange des contenus.}
\vspace{1em}
\begin{itemize}
    \item URL/URI (Tim Berners-Lee): principe d'adressage universel de toute information. Intégration des notions d'accès, partage ou appropriation des outils de publication.
    \item HTTP permet d'accéder à une ressource sans demander la permission au propriétaire du serveur.
    \item Tout média, toute information, toute donnée est susceptible sur le web d’être identifiée, extraite, associée, partagée, détournée.
\end{itemize}
    \textrightarrow{}  Culture de la libre circulation
\end{block}

Ces "règles" sont constamment négociées (par les grands groupes, les GAFAM, les chercheurs, la société civile, etc.)

\end{frame}

\begin{frame}{Exemple de travaux de recherche sur cette contre-culture}
    \centering
    {\small
    \textit{Les hippies ont projeté leur rêve d’exil et de refondation dans les échanges numériques et, pour cela, ils avaient besoin de couper les ponts avec un « réel » doublement décevant, en raison de la persistance de l’aliénation patriarcale et capitaliste, mais aussi de l’échec de la tentative de s’en émanciper en établissant dans ses marges des communautés contre-culturelles. Internet était un « ailleurs », le nouvel asile d’un projet d’émancipation avorté.}}\\
    \textbf{Fred Turner, \textit{Aux sources de l’utopie numérique}, C\&F édition, 2021 (2006)}
\end{frame}

\begin{frame}{Exemple de travaux de recherche sur cette contre-culture}
    \centering
    {\small \textit{Néanmoins, la rhétorique de « l’informationalisme en pair à pair », rejoignant en cela la rhétorique de la conscience dont elle était issue, occulte fortement les infrastructures matérielles et techniques dont dépendent à la fois l’Internet et les vies des membres de la génération numérique. Derrière le fantasme d’un flux d’information sans contrainte, s’étant la réalité de millions de claviers en plastique, de puces en silicium, de moniteurs aux écrans de verre, et de kilomètres de câbles à l’infini.}}\\
    \textbf{Fred Turner, \textit{Aux sources de l’utopie numérique}, C\&F édition, 2021 (2006)}
\end{frame}

\section{Comment ces utopies du Web ont-elles bouleversé nos protocoles éditoriaux ? Exemple avec Wikipédia}
\begin{frame}{Illustration avec un outil de tous les jours : Wikipédia}
Définie comme une encyclopédie collaborative, en ligne, à laquelle tout le monde peut donc gratuitement accéder et librement contribuer.\\
Revenir sur l'opposition fondamentale que certains établissent entre les modèles imprimés et le modèle numérique, afin de démontrer deux choses :
\begin{itemize}
    \item Général: le passage de l'imprimé au numérique n'est pas la révolution numérique que l'on veut faire croire.
    
    \item Cas particulier de l'encyclopédie: ce n'est pas la disparition d'un savoir faire, mais l'apparition d'un nouveau modèle épistémique, un modèle étroitement lié au fonctionnement même des media imprimés et numériques.
\end{itemize}

\end{frame}

\begin{frame}{D'\textit{Universalis} à \textit{Wikipédia} : deux modèles et protocoles éditoriaux numériques pour deux conceptions encyclopédiques opposées}
    
    \begin{columns}[T,onlytextwidth]
\column{0.50\linewidth}
\textbf{Modèle homothétique}
\centering
\includegraphics[width=.80\columnwidth]{05_Axe1_protocoles/images/universalis.png}
\begin{itemize}
    \item Acces fermé par une plateforme institutionnel restreinte
    \item Autorité de l'expert
    \item Stabilité et clôture du texte
\end{itemize}
\column{0.50\linewidth}
\textbf{Modèle éditorialisé}
\centering
\includegraphics[width=.80\columnwidth]{05_Axe1_protocoles/images/wikipedia-une-femme.png}
\begin{itemize}
    \item Écriture processuelle : ouverture dans le temps
    \item Écriture collective
    \item Accès ouvert et gratuit, ancrage dans le web
\end{itemize}
\end{columns}
\end{frame}

\begin{frame}{Le modèle Universalis : une culture de l'imprimé}
    \begin{columns}[T,onlytextwidth]
\column{0.50\linewidth}

\begin{itemize}
    \item Existe depuis 1966
    \item Publiée par la société d'édition Encyclopædia Universalis SA
    \item Née sur un support imprimé.
\end{itemize}
\column{0.50\linewidth}
\centering
\includegraphics[width=6cm]{05_Axe1_protocoles/images/universalis_livre.png}
\end{columns}

\end{frame}

\begin{frame}{Le modèle Universalis : une culture de l'imprimé}
\begin{itemize}
    \item Un éditeur en chef
    \item Un réseau "d'experts" (universitaires)
    \item Une chaîne éditoriale séquentielle (écriture -> correction -> publication)
    \item Des articles signés, relativement stables
\end{itemize}    
\end{frame}

\begin{frame}{Le modèle Universalis : une culture de l'imprimé}
\begin{center}
    \includegraphics[width=4cm]{05_Axe1_protocoles/images/universalis_cdrom.png}
    \hspace{2em}
    \includegraphics[width=7cm]{05_Axe1_protocoles/images/universalis_web.png}
\end{center}
\end{frame}

\begin{frame}{Deux acteurs importants}

\begin{block}{L'éditeur, figure centrale garante du respect du protocole}
    \begin{itemize}
        \item Garant de la fiabilité des contenus (gère les auteurs, vérifie les références)
        \item Garant de la mise en forme (gestion iconographique, cartes mentales)
    \end{itemize}
\end{block}
    
\begin{block}{L'auteur, guidé par le protocole éditorial}
    \begin{itemize}
        \item Des auteurs sollicités pour leur expertise (universitaires, professeurs)
        \item Des articles commandés (tâches d'écriture définies en amont dans un cahier des charges)
        \item Une attention portée à l'écriture (uniformisation du style, travail de synthèse et de vulgarisation)
    \end{itemize}
\end{block}
\end{frame}

\begin{frame}{Universalis en ligne : un modèle "homothétique"}
    \centering
    
\includegraphics[width=12cm]{05_Axe1_protocoles/images/universalis_ex.png}
\end{frame}

\begin{frame}{Un peu d'histoire : de Nupedia à Wikipédia, histoire d'un échec}

\begin{itemize}
    \item \textbf{Content Management System}: programme permettant de créer un site internet sans que les utilisateurs n'aient besoin de compétences en informatiques
    \item \textbf{Wiki} = logiciel permettant la création, la modification et l'illustration collaboratives de pages à l'intérieur d'un site web.
\end{itemize}

\end{frame}

\begin{frame}{Le projet Nupedia (2000 - 2003)}
    \begin{columns}
        \column{.30\linewidth}
        \begin{itemize}
            \item Fondateur: Jimmy Wales
            \item Éditeur en chef: Larry Sanger
        \end{itemize}
        \column{.60\linewidth}
        \includegraphics[width=.90\textwidth]{05_Axe1_protocoles/images/Nupedia_main_page.png}
    \end{columns}
\end{frame}

\begin{frame}{Nupedia : l'échec du protocole homothétique}
    \centering
    {\small
    \textit{Wikipédia n’est pas née radicale. Avant elle, Nupédia, la première encyclopédie en ligne lancée par Jimmy Wales au sein de sa société, Bomis, recouvrait une forme très traditionnelle. Un éditeur en chef, Larry Sanger, recruté à cet effet, surveillait et validait les contributions d’auteurs choisis pour leur expertise (Levrel, 2006). Le processus d’édition et de validation des articles était très rigoureux et s’appuyait sur un collège d’experts certifiés. La petite histoire raconte que Larry Sanger vérifiait que les contributeurs aient bien un doctorat (PhD), en leur demandant de lui envoyer leur diplôme par fax (Lih, 2009, p. 38)}}\\
    \vspace{1em}
    \textbf{Dominique Cardon, \textit{"Surveiller sans punir. La gouvernance de Wikipédia", Wikipédia, objet scientifique non identifié}, édité par Lionel Barbe et al., Presses universitaires de Paris Nanterre, 2015.}
\end{frame}

\begin{frame}{Nupedia, l'encyclopédie "incunable" ?}
    \centering
    {\small
    \textit{Au bout d’un an de fonctionnement, Nupédia ne parviendra pas à réunir plus de 20 articles : le processus était long, l’investissement des experts dilettante, le financement du projet fragile. Nupédia aurait très vite été oublié si, à titre expérimental, Larry Sanger et Jimmy Wales n’avaient décidé en janvier 2001 de lancer un wiki destiné à élaborer des ébauches d’articles pour Nupédia en permettant aux experts et à d’autres, d’écrire, de corriger et d’annoter les textes. À une vitesse stupéfiante, Wikipédia, simple « bac à sable » de Nupédia, va très vite se développer en suscitant l’attention d’internautes non experts, pendant que Nupédia sera définitivement déserté.}}\\
    \vspace{1em}
    \textbf{Dominique Cardon, \textit{"Surveiller sans punir. La gouvernance de Wikipédia", Wikipédia, objet scientifique non identifié}, édité par Lionel Barbe et al., Presses universitaires de Paris Nanterre, 2015.}

\end{frame}

\begin{frame}{De Nupedia à Wikipedia}
    \begin{block}{Une évolution du protocole éditorial...}
        \begin{itemize}
            \item Le CMS Wiki comme technologie de l'intellect
            \item Un modèle "horizontal" mais une gouvernance très stricte
            \item Un modèle extrêmement régulé (règles de contribution, de mise en forme, protocoles éditoriaux)
        \end{itemize}
    \end{block}

    \begin{block}{... et du concept de connaissance}
        \begin{itemize}
            \item Vers un principe de vérifiabilité
            \item Vers des communs de la connaissance
            \item Vers une autorité partagée
        \end{itemize}
    \end{block}
\end{frame}

\begin{frame}{Le CMS wiki : une interface régulatrice du collectif}
   Dominique Cardon (sociologue) explique que: le CMS (plateforme reposant sur un protocole éditorial co-construit entre les utilisateurs) suppose une régulation plus ou moins forte des comportements et des publications.\\
\begin{block}{}
\centering
     \textit{Les wikipédiens "ont choisi de confier à l’interface elle-même et aux règles de l’encyclopédie le soin de prendre en charge et de garantir que les comportements des wikipédiens soient bien orientées vers le savoir épistémique."}
\end{block}

\textrightarrow{} Le dispositif wiki fait agir les usagers en « wikipédiens » au profit du projet encyclopédique. 
\end{frame}

\begin{frame}{Le CMS wiki : une interface régulatrice du collectif}
\begin{itemize}
    \item Interactions avec les autres rédacteurs, les notifications apposées sur leurs contributions \textrightarrow{} les participants plient, adaptent et ajustent leurs comportements aux règles collectives.
    \item Ce sont les ajustements mutuels qui permettent aux personnes d'intéragir sur le wiki
    \item Interface de Wikipédia \textrightarrow{} contiguïté serrée entre l'espace des pratiques et celui des normes de la pratique.
\end{itemize} 
(Tiré de Dominique Cardon, \textit{"Surveiller sans punir. La gouvernance de Wikipédia", Wikipédia, objet scientifique non identifié}, édité par Lionel Barbe et al., Presses universitaires de Paris Nanterre, 2015.)
\end{frame}



\section{Le protocole éditorial du CMS Wiki}
\begin{frame}{Le protocole éditorial du CMS Wiki}
\begin{itemize}
    \item Un historique des modifications
    \item Une documentation des modifications
    \item Des espaces de débats ("Discussion")
    \item Des principes d'auto-évaluation et d'auto-régulation (labels, bots, etc.)
\end{itemize}

\end{frame}

\begin{frame}{Historique des modifications}
\centering
    \includegraphics[width=5cm]{05_Axe1_protocoles/images/wikipedia-une-femme-dif-historique.png}
    \includegraphics[width=5cm]{05_Axe1_protocoles/images/historique-une-femme.png}
\end{frame}

\begin{frame}{Documentation des modifications}
    \centering
    \includegraphics[width=.80\textwidth]{05_Axe1_protocoles/images/wikipedia-une-femme-commit.png}
\end{frame}

\begin{frame}{Espaces de débat}
    \centering
    \includegraphics[width=5cm]{05_Axe1_protocoles/images/wikipedia-une-femme-admissibilite-repose.png}
    \includegraphics[width=5cm]{05_Axe1_protocoles/images/wikipedia-une-femme-admissibilite.png}
\end{frame}

\begin{frame}{Système d'auto-évaluation et d'auto-régulation}
    \centering
    \includegraphics[width=.80\textwidth]{05_Axe1_protocoles/images/wikipedia-une-femme-auto-eval.png}
\end{frame}


\begin{frame}{Un savoir "instable", ou un savoir "éditorialisé" ?}
   \begin{itemize}
       \item Sur Wikipédia, une mise à jour constante
       \item Des contributions exponentielles
       \item Instabilité de l'encyclopédie, instabilité du savoir ?
   \end{itemize}


   \url{https://www.youtube.com/watch?v=ZOa8PAgg8ZI}


\end{frame}

\begin{frame}{Conclusion}
    La philosophie propre aux protocoles du web a participé à reconfigurer notre culture et notre manière de penser l'autorité. Ces protocoles ont encouragé l'émergence d'un nouveau modèle épistémique dont Wikipédia, incarnation de la nouvelle génération encyclopédique, peut témoigner. En passant du paradigme de l'édition à celui de l'éditorialisation, Wikipédia défend l'idée d'un savoir processuel, instable, sans cesse en construction. Un savoir qui, par ailleurs, doit se négocier au sein de d'une communauté - redéfinissant en cela notre conception même de l'autorité, sur laquelle s'était fondé le système éditorial moderne.
\end{frame}

\begin{frame}{Atelier: écrire, éditer sur un Wiki}
    \centering
    
\includegraphics[width=5cm]{05_Axe1_protocoles/images/qr-code-atelier-wiki.png}
\end{frame}
\end{document}